% =============================================================================
% SEÇÃO 4: RESULTADOS E DISCUSSÃO
% Diretrizes do Professor (Seção 3.2 do edital):
% - O clímax do trabalho.
% - Apresente gráficos de desempenho, tempo de execução ou tabelas de assertividade.
% - Discuta criticamente os resultados (mais do que apenas mostrar números):
%   * O que os resultados significam?
%   * Houve algum gargalo computacional?
%   * O algoritmo se comportou como esperado quando a entrada cresceu (complexidade)?
% =============================================================================

\section{Resultados e Discussão}
\label{sec:resultados}

% TODO: Escreva a análise de Resultados e Discussão aqui:
% 1. Apresentação das métricas topológicas e estatísticas computadas;
% 2. Tabelas comparativas de resultados e assertividade;
% 3. Gráficos de desempenho e visualizações das redes;
% 4. Discussão crítica aprofundada dos significados e anomalias encontradas;
% 5. Avaliação do comportamento assintótico e escalabilidade dos algoritmos.

% -----------------------------------------------------------------------------
% EXEMPLO DE INCLUSÃO DE TABELA:
% -----------------------------------------------------------------------------
% \begin{table}[!htbp]
% \centering
% \caption{Métricas estruturais das redes analisadas.}
% \label{tab:metricas}
% \begin{tabular}{lcc}
% \toprule
% \textbf{Métrica} & \textbf{Descrição} & \textbf{Valor} \\
% \midrule
% $|V|$ & Número de Vértices & -- \\
% $|E|$ & Número de Arestas & -- \\
% \bottomrule
% \end{tabular}
% \end{table}

% -----------------------------------------------------------------------------
% EXEMPLO DE INCLUSÃO DE GRÁFICO / FIGURA:
% -----------------------------------------------------------------------------
% \begin{figure}[!htbp]
% \centering
% \includegraphics[width=\columnwidth]{figures/grafico_exemplo.png}
% \caption{Descrição do gráfico e comportamento observado.}
% \label{fig:grafico}
% \end{figure}
